\documentclass{article}
\usepackage[T1]{fontenc}
\usepackage{lmodern}
\usepackage{graphicx}
\usepackage{amsmath,amssymb}
\usepackage[a4paper,margin=2cm]{geometry}
\usepackage[absolute,overlay]{textpos}
\setlength{\TPHorizModule}{1mm}
\setlength{\TPVertModule}{1mm}
\usepackage[indonesian]{babel}
\usepackage{hyperref}
\setlength{\emergencystretch}{2em}
% unit: Halaman 9

\begin{document}

% === Halaman 9 (python/native) ===

9

Enkripsi 

1. Preprocessing: Pesan yang akan dienkripsi dikodekan menjadi integer m. 

Jika m berukuran besar, pecah m menjadi blok-blok plainteks yang lebih 

kecil: m1, m2,  m3, …  ( syarat: 0  mi  < n – 1)

2. Hitung cipherteks c untuk plainteks m menggunakan  kunci publik e 

dengan persamaan 

                   c = me mod n

     Jika m dipecah menjadi m1, m2,  m3, … , enkripsi setiap mi menjadi

               ci = mi

e mod n

\end{document}
